\subsection{积分子群}

定义 1. 设 G 是李群。G 的一个积分子群是一个子群 H，带有连通李群结构，且 H 到 G 的规范嵌入是浸入。

\(G\) 的一参数子群是 \(G\) 的一个一维积分子群。

设 H 是 G 的积分子群，i 是 H 到 G 的规范嵌入。则 \(L(i)\) 确定了从 \(L(H)\) 到 \(L(G)\) 中一个具有拓扑补空间的李子代数的同构。以下将 \(L(H)\) 与其在 \(L(i)\) 下的像识别。

例。(1) \(\mathbf{G}\) 的一个连通李子群是 \(\mathbf{G}\) 的一个积分子群。

\begin{enumerate}
\def\labelenumi{(\arabic{enumi})}
\setcounter{enumi}{1}
\item
  假设 \(\mathbf{G}\) 是有限维的。令 \(\mathbf{H}\) 为 \(\mathbf{G}\) 的一个子群；赋予它由 \(\mathbf{G}\) 上的李群结构诱导的结构（§ 4, no. 5, Definition 3）。则其单位元分支 \(\mathrm{H}_0\) 是 \(\mathbf{G}\) 的一个积分子群，且 \(\mathbf{H}\) 在 \(e\) 处的切子代数是 \(\mathrm{L}(\mathrm{H}_0)\)（§ 4, no. 5, Proposition 9 (ii)）。
\item
  设 \(\mathbf{G}\) 是复李群，\(\mathbf{H}\) 是 \(\mathbf{G}\) 的积分子群，\(\mathbf{G}_1\)（分别为 \(\mathbf{H}_1\)）是 \(\mathbf{G}\)（分别为 \(\mathbf{H}\)）的底层实李群。则 \(\mathbf{H}_1\) 是 \(\mathbf{G}_1\) 的积分子群，且 \(\mathbf{L}(\mathbf{H}_1)\) 是 \(\mathbf{L}(\mathbf{H})\) 的底层实李代数。
\end{enumerate}

定理 2. 设 G 是李群。

\begin{enumerate}
\def\labelenumi{(\roman{enumi})}
\item
  映射 \(\mathbf{H} \mapsto \mathbf{L}(\mathbf{H})\) 把 \(\mathbf{G}\) 的积分子群集合双射到 \(\mathbf{L}(\mathbf{G})\) 中具有拓扑补空间的李子代数集合。
\item
  设 \(\mathbf{H}\) 是 \(\mathbf{G}\) 的积分子群。\(\mathbf{G}\) 的每个连通李子群芽，其李代数为 \(\mathbf{L}(\mathbf{H})\)，都是 \(\mathbf{H}\) 的开子流形，并生成 \(\mathbf{H}\)。
\end{enumerate}

\begin{enumerate}
\def\labelenumi{(\alph{enumi})}
\item
  令 \(\mathfrak{b}\) 为 \(\mathrm{L}(\mathrm{G})\) 中具有拓扑补空间的李子代数。令 \(\mathrm{H}_{1}\) 为 \(\mathrm{G}\) 的李子群芽，使得 \(\mathrm{L}(\mathrm{H}_{1}) = \mathfrak{b}\)（§ 4, Theorem 3）。可以选择 \(\mathrm{H}_{1}\) 使其连通。令 \(\mathrm{H}\) 为 \(\mathrm{G}\) 中由 \(\mathrm{H}_{1}\) 生成的子群。存在一个 \(\mathrm{H}\) 上的李群结构（§ 1, Corollary to Proposition 22），使得 \(\mathrm{H}_{1}\) 是 \(\mathrm{H}\) 的开子流形，且 \(\mathrm{H}\) 到 \(\mathrm{G}\) 的规范嵌入是浸入。由于 \(\mathrm{H}_{1}\) 连通，\(\mathrm{H}\) 连通，因而是 \(\mathrm{G}\) 的积分子群。于是 \(\mathrm{L}(\mathrm{H}) = \mathrm{L}(\mathrm{H}_{1}) = \mathfrak{b}\)。这证明了 (i) 中所考虑的映射是满射。
\item
  设 H 是 G 的积分子群，\(N_{1}\) 是李代数为 L(H) 的 G 的连通李子群芽。由于 H 到 G 的规范嵌入是浸入，存在 H 的一个开子群芽 \(H_{1}\)，它同时是 G 的子流形，因而是李代数为 L(H) 的 G 的李子群芽。另一方面，令 N 为由 \(N_{1}\) 生成的 G 的子群；由证明的 (a) 部分，N 带有 G 的积分子群结构，使 \(N_{1}\) 是 N 的开子流形。由 §4, Theorem 3，\(H_{1} \cap N_{1}\) 在 \(H_{1}\) 和 \(N_{1}\) 中都是开集。因此，由 \(H_{1} \cap N_{1}\) 生成的 G 的子群一方面等于 H，另一方面等于 N。所以李群 H 和 N 相等。这证明了 (ii)，也证明了 (i) 中所考虑的映射是单射。
\end{enumerate}

注记 1. 设 H 是 G 的积分子群。令 Y 为与 L(H) 相关的 G 的左叶层。若 \(g \in G\)，则通过 \(\gamma(g)\) 赋予 gH 由 H 的流形结构导出的流形结构。由 Differentiable and Analytic Manifolds, R, 9.3.2，gH 到 Y 的规范嵌入是态射。这个态射是 étale 的。因此，Y 的极大连通叶是模 H 的左陪集。

命题 1. 设 G 和 M 是李群，H 是 G 的积分子群，\(\phi\) 是从 M 到 G 的态射，且 \(\mathrm{L}(\phi)(\mathrm{L}(\mathrm{M})) \subset \mathrm{L}(\mathrm{H})\)。假设 M 连通。则 \(\phi(\mathrm{M}) \subset \mathrm{H}\)，并且把 \(\phi\) 看作从 M 到 H 的映射时，它是李群态射。

在注记 1 的记号下，\(\phi\) 是从 M 到 Y 的态射（Differentiable and Analytic Manifolds, R, 9.3.2），由于 M 连通，故 \(\phi(M) \subset H\)。

推论 1. 设 G 和 H 是李群，\(\phi\) 是从 G 到 H 的李群态射，N 是 \(\phi\) 的核，且 \(h = L(\phi)\)。假设 G 连通且 H 有限维。

\begin{enumerate}
\def\labelenumi{(\roman{enumi})}
\item
  N 是 G 的李子群，且 L(N) = Ker h。
\item
  令 \(\mathbf{H}'\) 为 \(\mathbf{H}\) 的积分子群，其李代数为 \(\operatorname{Im} h\)。则 \(\phi(\mathbf{G}) = \mathbf{H}'\)。
\item
  将 \(\phi\) 通过商群得到的从 G/N 到 H' 的映射是李群同构。
\end{enumerate}

\((i)\) 已经证明（§ 3, no. 8, Proposition 28）。

令 \(\psi\) 为将 \(\phi\) 通过商群得到的从 G/N 到 H 的李群态射；它是浸入（§3, no. 8, Proposition 28）。由命题 1，\(\psi\) 是从 G/N 到 H' 的李群态射。这个态射是 étale 的，由于 H' 连通，故 \(\psi\) (G/N) = H'；这证明了 (ii)。于是 \(\psi\colon G/N \to H'\) 是双射，也是李群同构，从而证明 (iii)。

推论 2. 设 \(\mathbf{G}\) 是李群，\(\mathrm{H}_1\) 和 \(\mathrm{H}_2\) 是 \(\mathbf{G}\) 的积分子群。若 \(\mathrm{L}(\mathrm{H}_2) \subset \mathrm{L}(\mathrm{H}_1)\)，则 \(\mathrm{H}_2\) 是 \(\mathrm{H}_1\) 的积分子群。

令 \(i_1: \mathrm{H}_1 \to \mathrm{G}\)、\(i_2: \mathrm{H}_2 \to \mathrm{G}\) 为规范嵌入。于是

\[
\mathrm{L} (i _ {2}) (\mathrm{L} (\mathrm{H} _ {2})) = \mathrm{L} (\mathrm{H} _ {2}) \subset \mathrm{L} (\mathrm{H} _ {1}).
\]

由命题 1，\(i_{2}\) 是从 \(H_{2}\) 到 \(H_{1}\) 的解析映射，甚至是从 \(H_{2}\) 到 \(H_{1}\) 的浸入；这是因为 \(\mathrm{L}(i_{2})\) 是从 \(\mathrm{L}(\mathrm{H}_{2})\) 到 \(\mathrm{L}(\mathrm{H}_{1})\) 的某个具有拓扑补空间的子代数的同构。

推论 3. 设 \(\mathbf{G}\) 是有限维李群，\((\mathrm{H}_i)_{i\in I}\) 是 \(\mathbf{G}\) 的一族李子群。则 \(\mathrm{H} = \bigcap_{i\in I}\mathrm{H}_i\) 是 \(\mathbf{G}\) 的李子群，并且

\[
\mathrm{L} (\mathrm{H}) = \bigcap_ {i \in \mathrm{I}} \mathrm{L} (\mathrm{H} _ {i}).
\]

存在 I 的一个有限子集 J，使得 \(\bigcap_{i\in J}L(H_i)\) 等于所有 \(\mathrm{L}(\mathrm{H}_i)\) 的交 M。我们知道 \(\mathrm{H}^* = \bigcap_{i\in J}\mathrm{H}_i\) 是满足 \(\mathrm{L}(\mathrm{H}^*) = \mathrm{M}\) 的李子群（§ 3, no. 8, Corollary 2 to Proposition 29）。令 \(\mathrm{H}_0\) 为 \(\mathrm{H}^*\) 的单位元分支。它是 G 的李子群，且 \(\mathrm{L}(\mathrm{H}_0) = \mathrm{M}\)。由推论 2，\(\mathrm{H}_0\subset \mathrm{H}_i\) 对所有 \(i\) 成立，因而 \(\mathrm{H}_0\subset \mathrm{H}\subset \mathrm{H}^*\)，从而得到推论。

推论 4. 设 G 是有限维连通李群。以下条件等价：

\begin{enumerate}
\def\labelenumi{(\roman{enumi})}
\item
  G 是幺模的（Integration, Chapter VII, § 1, no. 3, Definition 3）；
\item
  对所有 g ∈ G，都有 det Ad g = 1；
\item
  \(\operatorname{Tr} \operatorname{ad} a = 0\) 对所有 \(a \in \mathrm{L}(\mathrm{G})\) 都成立。
\end{enumerate}

映射 \(g \mapsto \operatorname{det} \operatorname{Ad} g\) 是一个态射 \(\phi\)，它从 \(G\) 到 \(K^*\)。由 § 3, Propositions 35 (no. 10) and 44 (no. 12)，有 \(L(\phi)a = Tr\operatorname{ad} a\) 对所有 \(a \in L(G)\) 成立。显然 \(\operatorname{Im} L(\phi) = \{0\}\) 或 \(K\)。在第一种（分别为第二种）情形，由推论 1，\(\operatorname{Im} \phi = \{1\}\)（分别为 \(\operatorname{Im} \phi = K^*\)），因而根据 § 3, no. 16, Corollary to Proposition 55，\(G\) 是幺模的（分别为非幺模的）。

命题 2. 设 \(\mathbf{G}\) 是有限维李群，\(\mathbf{H}\) 是 \(\mathbf{G}\) 的积分子群。以下条件等价：

\begin{enumerate}
\def\labelenumi{(\roman{enumi})}
\item
  H 是闭的；
\item
  H 上的拓扑由 G 上的拓扑诱导；
\item
  H 是 G 的李子群。
\end{enumerate}

\((i) \Rightarrow (iii)\)：这由 § 1, Propositions 2 (iv) (no. 1) and 14 (iii) (no. 7) 得出。

\((iii) \Rightarrow (ii)\)：显然。

\((ii) \Rightarrow (i)\)：若 H 上的拓扑由 G 上的拓扑诱导，则 H 是闭的，因为 H 是完备的（§ 1, no. 1, Proposition 1）。

命题 3. 设 G 是李群，H 是 G 的积分子群，M 是非空连通解析流形，f 是从 M 到 G 的映射，且 \(r \in N_{K}\)。考虑以下条件：

\begin{enumerate}
\def\labelenumi{(\roman{enumi})}
\item
  \(f\) 属于 \(\mathbf{C}^r\) 类，且 \(f(\mathbf{M})\subset \mathbf{H}\)；
\item
  \(f(\mathbf{M})\subset \mathbf{H}\)，并且把 \(f\) 看作从 \(\mathbf{M}\) 到 \(\mathbf{H}\) 的映射时，属于 \(\mathbf{C}^r\) 类
\item
  \(f\) 属于 \(\mathbf{C}^r\) 类，\(f(\mathbf{M})\) 与 \(\mathbf{H}\) 相交，并且 \(\mathrm{T}_m(\mathbf{M})\) 的像包含于 \(f(m)\) . \(\mathbf{L}(\mathbf{H})\)，对所有 \(m \in \mathbf{M}\)。
\end{enumerate}

则 (ii) \(\Leftrightarrow\) (iii) \(\Rightarrow\) (i)。若 H 上的拓扑有可数基，则三个条件等价。

\((ii) \Rightarrow (i)\) 以及 \((ii) \Rightarrow (iii)\)：显然。

\((iii) \Rightarrow (ii)\)：设条件 (iii) 成立。由 Differentiable and Analytic Manifolds, R, 9.2.8，\(f\) 是一个 \(\mathbf{C}^r\) 类态射，从 M 到与 \(\mathbf{L}(\mathbf{H})\) 相关的左叶层。由于 M 连通，\(f(\mathbf{M}) \subset \mathbf{H}\)。

若 H 上的拓扑有可数基，则条件 (i) 蕴含 f 是从 M 到 H 的 \(C^{r}\) 类映射（Differentiable and Analytic Manifolds, R, 9.2.8）；因此 (i) \(\Rightarrow\) (ii)。

推论 1. 设 G 是有限维李群，H 是 G 的积分子群。则 H 在 e 处的切李子代数（§ 4, no. 5, Definitions 2 and 3）是 L(H)，且 H 上的李群结构是由 G 上的结构诱导的结构。

由于 H 连通且有限维，其拓扑有可数基。

推论 2. 设 G 是李群，\(H_{1}\) 和 \(H_{2}\) 是 G 的积分子群。假设 \(H_{1}\) 上的拓扑有可数基。则

\[
\mathrm{H} _ {2} \subset \mathrm{H} _ {1} \Leftrightarrow \mathrm{L} (\mathrm{H} _ {2}) \subset \mathrm{L} (\mathrm{H} _ {1})
\]

并且，若这些条件成立，\(\mathbf{H}_2\) 是 \(\mathbf{H}_1\) 的积分子群。

最后一个断言以及蕴含 \(\mathrm{L}(\mathrm{H}_{2})\subset\mathrm{L}(\mathrm{H}_{1})\Rightarrow\mathrm{H}_{2}\subset\mathrm{H}_{1}\) 由 Proposition 1 的 Corollary 2 得出。逆向蕴含由命题 3 得出。

推论 3. 设 G 是李群，\(H_{1}\) 和 \(H_{2}\) 是 G 的积分子群，且它们的拓扑有可数基。若 \(H_{1}\) 和 \(H_{2}\) 具有相同的底层集合，则 \(H_{1}\) 和 \(H_{2}\) 上的李群结构相同。

这由推论 2 得出。

注记 2. 设 G 是有限维李群，H 是 G 的子群。若 H 上存在李群结构 S，使得带有 S 的 H 是 G 的积分子群，则我们约定称 H 是 G 的积分子群，这是一种语言上的滥用。由 Proposition 3 的 Corollary 3，若 S 存在，则 S 唯一。

注记 3. 设 V 是 \(\mathbf{C}^r\) 类流形。令 M 为 V 的子集，\(x\) 和 \(y\) 为 M 的元素。考虑以下性质：

\(P_{M,x,y}\) ：存在 I、\(x_{0}, x_{1}, \ldots, x_{n}, f_{1}, \ldots, f_{n}\)，使得：(a) I 是 K 的连通开子集；(b) \(x_{0}, \ldots, x_{n}\) 属于 M，\(x_{0} = x, x_{n} = y\)；(c) 对 \(1 \leqslant i \leqslant n\)，\(f_{i}\) 是从 I 到 V 的 \(C^{r}\) 类映射，取值为 \(x_{i-1}\) 和 \(x_{i}\)，且 \(f_{i}(I) \subset M\)。

称 M 为 V 的 \(C^{r}\)-连通子集，当且仅当对 M 中所有元素 x, y，性质 \(P_{M,x,y}\) 都成立。

命题 4. 设 G 是有限维李群，H 是 G 的子群。令 \(r \in \mathbf{N}_{\mathbb{K}}\)。以下条件等价：

\begin{enumerate}
\def\labelenumi{(\roman{enumi})}
\item
  H 是 G 的积分子群；
\item
  赋予 H 由 G 上的李群结构诱导的李群结构后，H 是连通的；
\item
  H 是 \(\mathbf{C}^r\)-连通的。
\end{enumerate}

\((ii) \Rightarrow (i)\)：显然。

\((i) \Rightarrow (iii)\)：设 H 带有李群结构，使 H 成为 G 的积分子群。使用注记 3 的记号，\(y \in H\) 中满足性质 \(\mathrm{P}_{\mathrm{H},e,y}\) 的元素所成的集合是 H 的一个开子群。由于 H 连通，这个子群等于 H，因而条件 (iii) 成立。

\((iii) \Rightarrow (ii)\)：设条件 (iii) 成立，并赋予 H 由 G 上的李群结构诱导的结构。令 \(\mathfrak{h}\) 为 H 在 \(e\) 处的切子代数。H 的单位元分支 \(\mathrm{H}_0\) 是 G 的积分子群，且 \(\mathrm{L}(\mathrm{H}_0) = \mathfrak{h}\)。我们证明 \(\mathrm{H} = \mathrm{H}_0\)。只需证明如下事实：令 I 为 K 的连通开子集，\(f\) 为从 I 到 G 的 \(\mathbf{C}^r\) 类映射，使 \(f(\mathrm{I}) \subset \mathrm{H}\)，令 \(\lambda\) 和 \(\mu\) 为 I 中两点；若 \(f(\lambda) \in \mathrm{H}_0\)，则 \(f(\mu) \in \mathrm{H}_0\)。但是，对所有 \(\nu \in \mathrm{I}\)，都有 \((\mathrm{T}_{\nu}f)(\mathrm{K}) \subset f(\nu)\mathfrak{h}\)，这是由 \(\mathfrak{h}\) 的定义得出的，所以我们的断言由命题 3 得出。

注记 4. 若 \(\mathbf{K} = \mathbf{R}\)，则 \(\mathbf{G}\) 的积分子群还可以刻画为这样的子群：赋予它们由 \(\mathbf{G}\) 的拓扑诱导的拓扑后，它们是弧连通的（§ 8, Exercise 4）。然而，子群可能是连通的，却不是积分子群（Commutative Algebra, Chapter VI, § 9, Exercise 2）。

推论。设 G 是有限维李群，\(H_{1}\) 和 \(H_{2}\) 是 G 的两个积分子群。由 \(H_{1}\) 和 \(H_{2}\) 生成的 G 的子群以及 G 的子群 ( \(H_{1}, H_{2}\) ) 都是 G 的积分子群。

G 的子群 (G, G) 不总是闭的（§ 9, Exercise 6）。

回忆（§ 3, no. 11, Corollary 5 to Proposition 41）：若 a 是有限维代数，则 \(\operatorname{Aut}(a)\) 是 \(\mathbf{GL}(a)\) 的李子群，且 \(\mathrm{L}(\mathrm{Aut}(a))\) 是 \(a\) 的导子李代数。

定义 2. 设 \(\mathfrak{a}\) 是有限维李代数。令 \(\operatorname{Ad}(\mathfrak{a})\) 或 \(\operatorname{Int}(\mathfrak{a})\) 表示 \(\operatorname{Aut}(\mathfrak{a})\) 中李代数为 \(\operatorname{ad}(\mathfrak{a})\) 的积分子群。这个群的元素称为 \(\mathfrak{a}\) 的内自同构。

由结构传递，\(\operatorname{ad}(a)\) 在 \(\operatorname{Aut}(a)\) 下不变，因而 \(\operatorname{Int}(a)\) 是 \(\operatorname{Aut}(a)\) 的正规子群。由 § 4, no. 4, Corollary 1 to Proposition 8 以及 \(\operatorname{Int}(a)\) 连通这一事实，\(\operatorname{Int}(a)\) 的元素是形如 \(\exp ad x\)（其中 \(x \in a\)）的自同构的有限乘积。一般而言，\(\operatorname{Int}(a)\) 不是 \(\operatorname{Aut}(a)\) 的李子群（练习 14）。

